\section{随机计算机的应用}

随机计算机是为解决自动控制领域的问题而开发的，其直接应用主要体现在自适应控制和自适应滤波领域. 用于过程辨识和在线优化的梯度技术是最简单的例子，这类强大的控制方法由于缺乏必要的硬件支持而未能得到充分应用. 使用传统计算机进行直接数字控制对于小型设备来说并不实用，而传统模拟计算机则因需要大量乘法器而成本高昂. 已有研究提出\(^{12}\)采用双电平或极性重合乘法\(^{10,11}\)作为一种利用数字集成电路廉价可靠地实现梯度技术的方法；然而，一项对最速下降计算中六种乘法技术的比较研究表明，在相同收敛时间下，极性重合乘法和继电器乘法产生的参数估计方差远大于等效的随机计算配置.\(^{5}\) 在这一特定计算中，随机乘法器可视为统计线性化\(^{13,14}\)的极性重合乘法器，而向输入信号添加锯齿波抖动以实现这种线性化的方法\(^{11,15,16}\)可以看作是一种获取伪随机序列的技术.

基于条件概率 Bayes 逆的最大似然预测是许多用于控制和模式识别的"学习机"的基础，但估计和预测的方程难以用传统计算元件实现，因此需要使用存储程序数字计算机来进行这项技术的实验. 然而，使用 ADDIE 的三输入变体，归一化似然比的估计和基于它们的预测都变得非常简单，只需很少的硬件即可实现.\(^{5}\)

随机计算在硬件经济性方面的理论基础在于 Rabin\(^{17}\) 和 Paz\(^{18}\) 的一个定理，该定理指出随机（或概率）自动机等价于一个通常具有更多状态的确定性自动机. 这一现象的一个直接实际例子可以在自适应阈值逻辑中找到，该逻辑被用于诸如 Perceptron\(^{19}\) 或 Adaline\(^{20}\) 之类的模式分类器中. 具有离散权重的自适应阈值逻辑元件在 Novikoff 条件\(^{21}\) 下不一定收敛，即使权重可以取到能实现线性可分的值，而等效的随机元件则可以在相同条件下收敛.\(^{5}\)

对这种差异的形象解释是：在具有连续权重的自适应阈值逻辑中，沿最速下降方向前进的路径在权重离散化后无法精确实现，此时存在多个"近似最速下降"的方向. 确定性逻辑必须"选择"其中一个方向，若选"错"了，可能会陷入错误决策的循环，而随机逻辑则以一定概率选择任意可能的下降方向，最终必定会选择正确的方向.
